\section{Design Decisions \& Rationale}

Engineering a production-quality, educational Computer Vision platform requires deliberate trade-offs across framework selection, inference execution, UI rendering, and state management. Table~\ref{tab:design_decisions} summarizes these architectural decisions alongside their concrete justifications.

\begin{table}[htbp]
    \centering
    \small
    \renewcommand{\arraystretch}{1.3}
    \begin{tabularx}{\textwidth}{|p{3.2cm}|p{3.0cm}|X|}
        \hline
        \textbf{Decision Area} & \textbf{Chosen Solution} & \textbf{Architectural Rationale \& Trade-Offs} \\
        \hline
        \textbf{Backend Framework} & FastAPI (Python 3.13) with Uvicorn ASGI & \textbf{Rationale:} Chosen over Flask or Django due to native asynchronous concurrency (`\texttt{async/await}`), automated OpenAPI/Swagger generation, and high throughput. Handles heavy multi-part video uploads without blocking event loops. \\
        \hline
        \textbf{Frontend Architecture} & Vanilla HTML5, Modern CSS3, ES6 JavaScript & \textbf{Rationale:} Deliberately avoids bulky frontend frameworks (React, Vue, Next.js). Eliminates multi-hundred-megabyte `\texttt{node\_modules}` dependencies, complex build configurations, and bundling overhead. Runs natively in any web browser via a simple Python HTTP daemon. \\
        \hline
        \textbf{Deep Learning Detector} & Ultralytics YOLOv8n (Nano Variant) & \textbf{Rationale:} Offers state-of-the-art mean Average Precision (mAP) with an anchor-free single-shot detection head. Weighs only $\sim6$\,MB, allowing instant automatic downloading on first use, and sustains real-time inference on consumer CPUs without requiring dedicated GPUs. \\
        \hline
        \textbf{Multi-Object Tracker} & ByteTrack Association & \textbf{Rationale:} Chosen over DeepSORT or SORT because ByteTrack associates both high-confidence and low-confidence detection boxes using a Kalman filter and Hungarian bipartite matching. Preserves track IDs through severe crowd occlusions with zero re-identification network overhead. \\
        \hline
        \textbf{Computer Vision Engine} & OpenCV Headless with NumPy 2.0 & \textbf{Rationale:} `\texttt{opencv-python-headless}` prevents GUI/X11 runtime crashes on Linux and headless environments while retaining full access to optimized C++ computer vision kernels. NumPy 2.0 provides SIMD vectorized matrix arithmetic. \\
        \hline
        \textbf{Visualization Engine} & Chart.js 4.4 \& HTML5 Canvas 2D & \textbf{Rationale:} Renders high-performance histograms, polar motion vector charts, and bounding box coordinates directly on native HTML5 Canvas elements with minimal memory footprint. \\
        \hline
        \textbf{UI Design System} & Dark Glassmorphism with High-Contrast Tokens & \textbf{Rationale:} Solves visual fatigue during prolonged laboratory sessions. High-contrast typography (`\texttt{\#f8fafc}` primary, `\texttt{\#cbd5e1}` secondary) and semi-opaque panels (`$86\%$` opacity) guarantee complete readability against dynamic ambient backgrounds. \\
        \hline
    \end{tabularx}
    \caption{Key Architectural Design Decisions and Rationale.}
    \label{tab:design_decisions}
\end{table}

\subsection{In-Depth Architectural Justifications}

\subsubsection{Why Decoupled Client-Server Over a Monolithic Desktop GUI?}
Desktop GUI toolkits such as Tkinter, PyQt, or wxPython are frequently used in simple student projects. However, they suffer from critical drawbacks: platform-specific UI rendering bugs, heavy installation requirements (e.g., Qt runtime libraries), and a complete lack of web accessibility. By implementing a decoupled web-based architecture (FastAPI backend + Vanilla web client), VisionLab achieves complete cross-platform portability: any device with a modern web browser can interact with the system without installing native desktop bindings.

\subsubsection{Why Stateless REST Over WebSockets for Core Processing?}
While WebSockets provide persistent bi-directional duplex channels, classical Computer Vision modules (filtering, edge extraction, clustering, single-shot detection) are inherently request-response transactions. Employing HTTP POST multipart/form-data requests ensures:
\begin{itemize}
    \item Strict adherence to RESTful principles and idempotent caching where applicable.
    \item Complete compatibility with standard benchmarking tools (`curl`, Postman, Swagger).
    \item Simplified error handling, status code semantics ($200$, $422$, $500$), and independent microservice scalability.
\end{itemize}
